\documentclass[11pt]{article}
\usepackage[utf8]{inputenc}
\usepackage[T1]{fontenc}
\usepackage{amsmath,amssymb,amsthm}
\usepackage[margin=1in]{geometry}
\usepackage{booktabs}
\usepackage[round]{natbib}
\setcitestyle{aysep={},yysep={;}}
\usepackage{hyperref}
\usepackage{microtype}
\emergencystretch=2em

\newtheorem{theorem}{Theorem}
\newtheorem{proposition}{Proposition}
\newtheorem{lemma}{Lemma}
\newtheorem{corollary}{Corollary}
\theoremstyle{remark}
\newtheorem{remark}{Remark}
\newcommand{\E}{\mathbb{E}}
\newcommand{\cD}{\mathcal{D}}
\newcommand{\cA}{\mathcal{A}}
\newcommand{\cP}{\mathcal{P}}
\newcommand{\Lat}{\mathrm{Lat}}
\newcommand{\Ind}{\mathrm{Ind}}
\newcommand{\pb}{\bar p}
\newcommand{\qb}{\bar q}

\title{Larger Crowds Do Not Restore Truthfulness in Peer Prediction with Heterogeneous Priors:\\
A Counterexample to Theorem~2 of Qiu, Carroll and Allen (ICLR 2026), with Machine-Checked Inequalities}
\author{Alejandro Panera Pascual (GalexY)\\ \small Universidad de Salamanca, Spain\\ \small \url{https://github.com/GalexY3}}
\date{November 2026 (version 1.0.1)}

\begin{document}
\maketitle

\begin{abstract}
\citet{qiu2026truthfulness} propose peer prediction for evaluating and training language models without ground truth: a participant's answer is scored by how much it helps a crowd of expert models predict the answers of the other participants. Their Theorem~2 states that, when the priors of participants and experts are drawn independently from a common distribution $\cD$ and the numbers of experts $m$ and participants $n$ are large, honest reporting is, with high probability, an $\varepsilon$-Bayesian Nash equilibrium. We show that the statement is false under the payoff normalization used in the steps of its proof, and that under the normalization printed in two of the proof's displayed equations it holds for every $\cD$ for a trivial reason, because every incentive then shrinks as $1/m$; in neither case does the proof establish it. In a two-type population with binary answers, where 90\% of priors follow a latent truth observed with noise and 10\% treat answers as independent and mostly~0, a participant with the minority prior gains at least $c_0=\tfrac15\log\tfrac{9}{41}+\tfrac45\log\tfrac{51}{19}\approx0.487$ nats per target by reporting~0 when its answer is~1, and at least $c_0/5\approx0.097$ nats per target averaged over its own answer, for every fraction of majority-type experts in $[4/5,1]$. The gain does not depend on $m$ or $n$, the hypotheses of the theorem, including its Assumption~1, are satisfied, and with probability tending to one some participant has the minority prior; so honest reporting is not an $\varepsilon$-equilibrium for any $\varepsilon<0.097$ once $m$ and $n$ exceed the theorem's own threshold. The profitable deviation, always answering~0, is uninformative. The counterexample survives the other readings of the equilibrium notion that we consider, including one in which deviations may not depend on the participant's own prior, refuted with non-exchangeable priors. The decisive gap in the proof is the replacement of the deviating participant's own prior by an aggregate of the experts' priors, which the passage to finite $n$ does not remove; the passage between arithmetic and geometric averages of expert probabilities, attributed to Hoeffding's inequality, is a second and independent gap. The key inequalities are machine-checked in Lean~4 with Mathlib, with statements fixed before the proofs were written, and the model is checked by exact computation.
\end{abstract}

\section{Introduction}

Peer prediction rewards an agent for reports that help predict the reports of others, and with a common prior it makes honest reporting an equilibrium without access to ground truth \citep{miller2005peer,schoenebeck2022two}. \citet{qiu2026truthfulness} bring this idea to the evaluation and post-training of large language models. Participants (models under evaluation) answer a question; experts (possibly much weaker models) report how likely each participant's answer is, alone and conditional on another participant's answer; a participant acting as \emph{source} is scored by how much its answer improves the experts' log-probability of a \emph{target} participant's answer. Their Theorem~1 restates the common-prior guarantee. Their Theorem~2, presented as novel and of independent interest, lifts it to heterogeneous priors: if the priors of participants and experts are drawn independently from one distribution $\cD$, then for large $m$ and $n$ honest reporting is, with probability $1-\delta$, an $\varepsilon$-Bayesian Nash equilibrium. The paper reads this as saying that making the pool of agents ``large and diverse'' is sufficient for incentive compatibility; after Theorem~1 it also states that the mechanism penalizes both deceptive and uninformative answers \citep[Sec.~3.3]{qiu2026truthfulness}.

The difficulty is that a large crowd of experts makes their \emph{average} prediction converge to the population average, not to the belief of any single participant. A participant whose prior differs from the population average maximizes its expected log score by choosing the report whose induced crowd prediction is closest, in Kullback--Leibler divergence, to its own belief about the target. That report need not be its answer, and more experts only stabilize the incentive to misreport. This is a Goodhart effect \citep{goodhart1975problems,strathern1997improving} in a precise sense: the score measures agreement with the crowd's prediction, and optimizing it leads a participant to stop conveying its information.

\paragraph{Contributions.}
\begin{itemize}
\item A counterexample to Theorem~2 under the reading used in its proof, in which each participant knows its own prior (Section~\ref{sec:caseA}). The deviation gains at least $0.487$ nats per target conditional on the participant's answer and $0.097$ averaged over it, uniformly in $m$ and $n$; the same population also refutes the variant in which the ex ante payoff averages over the draw of the participant's own prior.
\item A counterexample under the weaker, ex ante reading in which deviations cannot depend on the participant's own prior, using non-exchangeable priors (Section~\ref{sec:caseB}).
\item A proof that, with the normalization printed in equations (17) and (22) of the paper, the statement is trivially true (Lemma~\ref{lem:triv}), so that Theorem~2 is either false or trivial.
\item The location of the gaps in the proof in Appendix~C.2 of the paper; the decisive one is the replacement of the deviating participant's prior by the experts' aggregate (Section~\ref{sec:proof}).
\item Lean~4 proofs of the key inequalities with preregistered statements, and exact brute-force checks (Section~\ref{sec:lean}).
\end{itemize}
We checked the published version (OpenReview, ``Published as a conference paper at ICLR 2026'') against arXiv:2601.20299v1: the statement of Theorem~2 and the steps of its proof discussed here are the same in both.

\section{Setting}\label{sec:setting}

There are $n$ participants with answers $A_1,\dots,A_n\in\cA$ and $m$ experts. Expert $j$ reports probabilities $\Pr_j(A_t)$ and $\Pr_j(A_t\mid A_s)$. In Algorithm~2 of \citet{qiu2026truthfulness}, which is the version covered by Theorem~2, expert probabilities are averaged before the logarithm is taken:
\begin{gather}
S_s=\frac1{n-1}\sum_{t\neq s}\Big[\log\pb(A_t\mid A_s)-\log\qb(A_t)\Big],\label{eq:score}\\
\pb(y\mid x)=\frac1m\sum_j\Pr_j(A_t{=}y\mid A_s{=}x),\qquad \qb(y)=\frac1m\sum_j\Pr_j(A_t{=}y).\notag
\end{gather}
Expert weights are uniform, as in Theorem~2. Honest experts report the conditionals of their own priors. Under the honest profile, targets answer honestly.

\paragraph{Normalization.} Algorithm~2 returns the sum over targets of the bracket in \eqref{eq:score}. The proof of Theorem~2 divides it by $(n-1)m$ in its equations (17), (18) and (22), and by $n-1$ in (19)--(21). We use $1/(n-1)$, the per-target average in \eqref{eq:score}; with the plain sum, every gain below is multiplied by $n-1$. The factor $1/m$ cannot be meant literally:

\begin{lemma}[The literal normalization makes Theorem~2 trivial]\label{lem:triv}
Under Assumption~1 (below), $|\log\pb(y\mid x)-\log\qb(y)|\le I_0$ for all answers $x,y$ and all profiles of honest experts with priors in the support of $\cD$. Hence, with payoff $S_s/m$, no deviation of a participant changes its expected payoff by more than $2I_0/m$, and when $m$ satisfies the threshold \eqref{eq:thr} this is less than $\varepsilon$, for every $\cD$, $\varepsilon>0$ and $\delta\in(0,1]$.
\end{lemma}
\begin{proof}
$\pb(y\mid x)/\qb(y)=\sum_j\Pr_j(y)\,e^{\mathrm{pmi}_j(x;y)}\big/\sum_j\Pr_j(y)$ is a weighted average of $e^{\mathrm{pmi}_j(x;y)}$, and $|\mathrm{pmi}_j|\le I_0$. So every target term of \eqref{eq:score} lies in $[-I_0,I_0]$, as does $S_s$, and a deviation changes $S_s$ by at most $2I_0$. Since $|\cA|/\delta\ge2$, \eqref{eq:thr} gives $m>16\log 2\,(I_0+L_0)/\varepsilon>11I_0/\varepsilon$, so $2I_0/m<\varepsilon$.
\end{proof}

Under that reading Theorem~2 holds for a reason unrelated to crowds or priors, and the incentive to answer honestly vanishes as $m$ grows; the proof itself drops the factor at its step (19). We therefore take the payoff to be \eqref{eq:score}, and our conclusion is that Theorem~2 is false with the normalization its proof uses and trivial with the one printed in (17) and (22).

\begin{lemma}[Best response of a source]\label{lem:br}
Fix the expert reports and the honest behaviour of the targets. A source $s$ whose answer is $a$ and whose belief about a target's answer is $P_s(\cdot\mid a)$ obtains, from target $t$, expected score $\sum_y P_s(y\mid a)\log\pb(y\mid r)$ plus a term that does not depend on its report $r$. It therefore prefers the report $r$ that minimizes $\mathrm{KL}\big(P_s(\cdot\mid a)\,\|\,\pb(\cdot\mid r)\big)$. For binary answers and $b=P_s(A_t{=}1\mid a)$, the report-dependent part is
\begin{equation}
\mathrm{sc}(b,p)=b\log p+(1-b)\log(1-p),\qquad p=\pb(1\mid r).\label{eq:sc}
\end{equation}
\end{lemma}
\begin{proof}
The term $\log\qb(A_t)$ in \eqref{eq:score} does not involve $r$, and $\sum_yP\log\pb=-\mathrm{KL}(P\|\pb)-H(P)$.
\end{proof}

Honesty ($r=a$) is a best response only if $\pb(\cdot\mid a)$ is the closest of the induced crowd predictions to the source's own belief. With a common prior this holds, because the experts' prediction \emph{is} the source's belief. With heterogeneous priors it need not.

\paragraph{The statement.} Theorem~2 assumes that the priors $\cP_i^{\mathrm A}$ of participants and $\cP_j^{\mathrm J}$ of experts are drawn independently from a distribution $\cD$ over priors, and Assumption~1: the pointwise mutual information between the answers of two participants, within any prior in the support, is bounded by $I_0$, and log-ratios of marginals across priors are bounded by $L_0$.\footnote{The paper describes each prior, a distribution over $\cA^n$, as a vector in $[0,1]^{n|\cA|}$. That dimension fits only product priors, under which every pointwise mutual information is zero and the score \eqref{eq:score} is identically zero; we read priors as joint distributions over $\cA^n$, as Assumption~1 and the proof require. Similarly, the supremum defining $I_0$ is read over pairs of distinct participants, as the paper's Remark~3 describes it: with $i=j$ and $\hat a_i\neq\hat a_j$ it would be infinite for every $\cD$. With the diagonal terms $\hat a_i=\hat a_j$ included, $I_0$ becomes $\log5$ in case~A and the threshold $2791$ below becomes $3754$; the counterexample is unaffected.} It concludes that if
\begin{equation}
m,n\ \ge\ \frac{16(I_0+L_0)}{\varepsilon}\log\Big(\frac{I_0+L_0}{\varepsilon}+\frac{|\cA|}{\delta}\Big),\label{eq:thr}
\end{equation}
then with probability $1-\delta$ the honest profile is, ``ex ante (when $\cD$ and the instantiation of all $\cP_j^{\mathrm J}$ are known by the agents)'', an $\varepsilon$-Bayesian Nash equilibrium. We consider two readings:
\begin{itemize}
\item \textbf{Reading I.} Each participant knows its own prior; deviations may depend on it. This is the notation of the proof, which takes expectations under $\cP_s^{\mathrm A}$.
\item \textbf{Reading I$'$.} As Reading~I, but the participant's own prior is part of its type and the ex ante payoff also averages over the draw of that prior from $\cD$.
\item \textbf{Reading II.} Participants do not condition on their own prior: a deviation is a map $\sigma:\cA\to\cA$ shared by all types, evaluated ex ante over $\cD$.
\end{itemize}
Within Reading~I, the $\varepsilon$-equilibrium can be taken interim (conditional on the participant's own answer) or averaged over its own answer; we give the bound for both.

\section{Counterexample under Reading I}\label{sec:caseA}

Let $\cA=\{0,1\}$ and $\cD_A=0.9\,\delta_{\Lat}+0.1\,\delta_{\Ind}$, with two exchangeable priors defined for every $n$:
\begin{itemize}
\item $\Lat$ (latent truth): $Z\sim\mathrm{Bern}(1/2)$, and each $A_i=Z$ with probability $9/10$, independently given $Z$. For every pair, $P(A_t{=}1\mid A_s{=}1)=41/50$ and $P(A_t{=}1\mid A_s{=}0)=9/50$.
\item $\Ind$ (independent): the $A_i$ are i.i.d.\ $\mathrm{Bern}(1/5)$.
\end{itemize}
Assumption~1 holds with $I_0=\log(25/9)\approx1.02$ and $L_0=\log(5/2)\approx0.92$.

If a fraction $f$ of the experts has prior $\Lat$, \eqref{eq:score} gives
\begin{equation}
\pb(1\mid1)=f\tfrac{41}{50}+(1-f)\tfrac15=\tfrac{10+31f}{50},\qquad \pb(1\mid0)=f\tfrac{9}{50}+(1-f)\tfrac15=\tfrac{10-f}{50}.
\end{equation}
A participant with prior $\Ind$ believes $P(A_t{=}1)=1/5$ whatever its own answer, so $b=1/5$ in \eqref{eq:sc}.

\begin{proposition}[Lean: \texttt{caseA\_gain}, \texttt{caseA\_c0\_pos}]\label{prop:A}
For every $f\in[4/5,1]$,
\[
\mathrm{sc}\big(\tfrac15,\pb(1\mid0)\big)-\mathrm{sc}\big(\tfrac15,\pb(1\mid1)\big)\ \ge\ c_0=\tfrac15\log\tfrac{9}{41}+\tfrac45\log\tfrac{51}{19}\approx0.4866>0.
\]
\end{proposition}

So a participant with prior $\Ind$ whose answer is~1 gains at least $c_0$ per target by reporting~0; when its answer is~0, reporting~0 is honest. The strategy ``always~0'' thus beats honesty by at least $c_0$ conditional on answer~1, and by at least $c_0/5\approx0.097$ averaged over its own answer, since its prior gives answer~1 probability $1/5$. The gain is $0.524$ at $f=4/5$, $0.689$ at $f=9/10$ and $0.910$ at $f=1$, conditional on answer~1. It is the same for every target, so the average over targets in \eqref{eq:score} has the same gain, and it does not depend on $m$ or $n$.

\begin{corollary}\label{cor:A}
Under $\cD_A$, let $E$ be the event that the fraction of $\Lat$ experts lies in $[4/5,1]$ and some participant has prior $\Ind$. On $E$, the honest profile is not an $\varepsilon$-Bayesian Nash equilibrium for any $\varepsilon<c_0/5$ (averaged over the participant's answer), nor for any $\varepsilon<c_0$ (interim). Since $P(E)=(1-0.9^n)\big(1-P(\text{fraction}\notin[4/5,1])\big)\to1$, the conclusion of Theorem~2 fails for every $\delta<P(E)$ once $m$ and $n$ are large.
\end{corollary}

Table~\ref{tab:probs} gives exact probabilities. For $\delta=0.05$, the threshold \eqref{eq:thr} is $m,n\ge2791$ for $\varepsilon=c_0/10\approx0.049$, which refutes the averaged reading, and $m,n\ge494$ for $\varepsilon=c_0/2$, which refutes only the interim reading (the averaged gain never exceeds $0.182$). At those sizes $P(E)$ is $1$ to many decimal places.

Under Reading~I$'$ the minority type has probability $1/10$, so the gain is at least $c_0/50\approx0.0097$, still independent of $m$ and $n$ (threshold $17\,652$ for $\varepsilon=0.99\,c_0/50$ and $\delta=0.05$). If the agents do not know the fraction $f$ and average over the draw of the experts, the gain is at least $c_0\cdot P(f\in[4/5,1])$, because the conditional gain is non-negative for every $f\in[0,1]$.

\begin{table}[h]
\centering\small
\begin{tabular}{lccccc}
\toprule
$m$ & 50 & 100 & 200 & 500 & 1000\\
$P(\text{fraction }\Lat\notin[4/5,1])$ &$9.4\cdot10^{-3}$ & $8.1\cdot10^{-4}$ & $7.2\cdot10^{-6}$ & $7.9\cdot10^{-12}$ & $1.3\cdot10^{-21}$\\
\midrule
$n$ & 5 & 10 & 20 & 50 & 100\\
$P(\text{some }\Ind\text{ participant})=1-0.9^n$ &0.410 & 0.651 & 0.878 & 0.995 & 0.99997\\
\bottomrule
\end{tabular}
\caption{Exact probabilities under $\cD_A$ (binomial sums with rational arithmetic).}\label{tab:probs}
\end{table}

\section{Counterexample under Reading II}\label{sec:caseB}

Under Reading~II the participant cannot use its own prior, so its belief is the mixture over $\cD$. With exchangeable latent-class priors we did not find a counterexample in an exploratory search (this is not a proof that none exists). The hypothesis of Theorem~2 allows any $\cD$ over priors on $\cA^n$, however, and non-exchangeable priors suffice.

Take two ``star'' priors centred on participant~1: $A_1\sim\mathrm{Bern}(a)$ and the others i.i.d.\ given $A_1$ with $P(A_t{=}1\mid A_1{=}x)=r(x)$. Let $P_1=(a,r(1),r(0))=(1/20,9/20,1/20)$, $P_2=(19/20,1/5,19/20)$ and $\cD_B=\tfrac{7}{20}\delta_{P_1}+\tfrac{13}{20}\delta_{P_2}$. Assumption~1 holds with $I_0\approx2.72$ and $L_0\approx2.94$. Under the mixture, $P(A_1{=}0,A_t{=}1)=19/400$ and $P(A_1{=}0,A_t{=}0)=127/400$ (Lean: \texttt{caseB\_mixture}). With a fraction $f$ of experts with prior $P_1$, $\pb(1\mid1)=f\tfrac{9}{20}+(1-f)\tfrac15$ and $\pb(1\mid0)=f\tfrac1{20}+(1-f)\tfrac{19}{20}$.

\begin{proposition}[Lean: \texttt{caseB\_gain}, \texttt{caseB\_c1\_pos}]\label{prop:B}
For every $f\in[3/10,2/5]$, the ex ante gain of ``always~1'' over honesty for participant~1,
$\tfrac{146}{400}\big[\mathrm{sc}(\tfrac{19}{146},\pb(1\mid1))-\mathrm{sc}(\tfrac{19}{146},\pb(1\mid0))\big]$, is at least
$c_1=\big(19\log\tfrac{55}{136}+127\log\tfrac{70}{41}\big)/400\approx0.127>0$.
\end{proposition}

The gain is $0.175$ at $f=7/20$. The fraction of $P_1$ experts concentrates at $7/20$ and falls outside $[3/10,2/5]$ with probability $0.017$ at $m=500$ and $2.5\cdot10^{-6}$ at $m=2000$. The mechanism is that experts average \emph{conditionals}, whereas the participant's ex ante belief is the conditional of the \emph{mixture}; with different marginals the two differ.

\section{Where the proof fails}\label{sec:proof}

Appendix~C.2 of \citet{qiu2026truthfulness} first treats $n\to\infty$ under the assumption $\cP_i^{\mathrm A}=\bar\cP$ for all $i$, where $\log\bar\cP:=\E_{\cP\sim\cD}[\log\cP]$, and then removes that assumption for finite $n$. Its chain is (17) $=$ (18) $\ge$ (19) $=$ (20) $\ge$ (21) $\ge$ (22): steps (17)--(19) take expectations under the source's own prior $\cP_s^{\mathrm A}$, steps (20)--(21) under $\bar\cP$, and (22) returns to $\cP_s^{\mathrm A}$.

\paragraph{The decisive gap: the source's prior is replaced by $\bar\cP$.} The equality (19) $=$ (20) holds only if $\cP_s^{\mathrm A}=\bar\cP$, and step (21), which uses the non-negativity of the Kullback--Leibler divergence, requires that the distribution under which answers are drawn be the one that predicts them. For a source whose prior differs from $\bar\cP$ neither holds. Case~A shows that this is the operative failure, independently of any averaging: at $f=1$ all experts share the prior $\Lat$, so $\bar\cP=\Lat$ and the experts' average is exact, yet the $\Ind$ participant gains $0.910$ per target conditional on answer~1 ($0.182$ averaged over its answer).

\paragraph{Finite $n$.} The proof removes the assumption $\cP_i^{\mathrm A}=\bar\cP$ by noting that ``(20) is the mean of (17) over the choice of $\{\cP_i^{\mathrm A}\}$''. Averaging (17) over the source's prior drawn from $\cD$ gives an expectation under the arithmetic mixture $\E_\cD[\cP]$, not under $\bar\cP$, and it compares honesty only with deviations that do not depend on the source's prior. It does not cover Reading~I, and Section~\ref{sec:caseB} shows that the averaged statement also fails without exchangeability: there the experts average conditionals, while the participant's belief is the conditional of the mixture.

\paragraph{The Hoeffding step.} Steps (18) $\to$ (19) and (21) $\to$ (22) pass between the arithmetic average $\frac1m\sum_j\Pr_j$ computed by Algorithm~2 and the geometric mean $\bar\cP$, ``from Hoeffding's inequality''. Hoeffding's inequality \citep{hoeffding1963} concentrates the average around $\E_\cD[\Pr]$, an \emph{arithmetic} mean; the logarithm of the arithmetic mean differs from the mean of the logarithms by a Jensen gap that does not shrink with $m$ and can have either sign. With $\bar\cP$ computed on pairs of answers (the normalized geometric mean of the pairwise joint distributions) and $f=9/10$ in case~A, the deviation ``always~0'' of the $\Ind$ participant scores $0.1574$ under the arithmetic average and $0.1315$ under $\bar\cP$, so the passage (21) $\to$ (22) fails for every $\varepsilon<0.052$; and a source with prior $\Lat$ answering honestly scores $0.2180$ under the arithmetic average and $0.2245$ under $\bar\cP$, so (19) fails for every $\varepsilon<0.013$. For the paper's $\bar\cP$ on $\cA^n$ these values depend on $n$; in case~A, $\bar\cP$ approaches i.i.d.\ $\mathrm{Bern}(0.1075)$ answers as $n$ grows. This gap is real but secondary: the counterexample does not rely on it.

\paragraph{Algorithm~1.} Theorem~2 is stated for Algorithm~2. In Algorithm~1, which averages logarithms over experts and also scores the round $t=s$, the same deviation still gains, conditional on answer~1, $\big[(n-1)\,f\cdot\tfrac35\log\tfrac{41}{9}-(1-f)\log4\big]/n>0$ per round for $f\in[4/5,1]$ and $n\ge2$; the first term is $f\big(\tfrac45\log\tfrac{41}{9}+\tfrac15\log\tfrac{9}{41}\big)$, whose positivity is \texttt{alg1\_caseA\_pos} in Lean.

\section{Machine-checked statements}\label{sec:lean}

The file \texttt{PeerPredictionGD01.lean} \citep{demoura2021lean4,mathlib2020} proves the six statements named above. They are stated over the reals (\texttt{caseB\_mixture} over the rationals), with $\mathrm{sc}$ and $\pb$ as defined here. The statements were fixed in a preregistration, with SHA-256 freeze, before the proofs were written, and an automatic comparison found the nine declarations identical to the preregistered text. All six compile without \texttt{sorry}; \texttt{\#print axioms} reports only \texttt{propext}, \texttt{Classical.choice} and \texttt{Quot.sound}.

Lean certifies the scalar inequalities only. The model is checked by exact computation: Lemma~\ref{lem:br}, the conditionals of $\Lat$, the reduction of the $n$-participant score to per-pair terms (brute force over all answer profiles in $\cA^n$ for $n=2,3,4$ with rational arithmetic, matching the closed forms to within $10^{-16}$, together with the linearity of \eqref{eq:score} in the targets), the mixture belief of case~B, the constants of Assumption~1 and the binomial probabilities. An independent re-encoding written from this note, with Assumption~1 checked over all pairs and the gains for $n\le5$, reproduces every number.

\paragraph{Changes from version 1.0.0.} Version 1.0.1 follows an adversarial audit by two independent readers (fidelity to the paper; arithmetic and Lean). It adds the normalization paragraph and Lemma~\ref{lem:triv}, the footnote on the reading of the hypotheses, Reading~I$'$ and the remark on unknown $f$, and states which threshold refutes which reading. It rewrites Section~\ref{sec:proof}: version 1.0.0 illustrated step (19) with numbers that satisfy that step, and gave the Hoeffding step as the first gap. It also completes the Algorithm~1 expression and corrects the attribution of the sentence on uninformative answers. The counterexamples and every number in them are unchanged.

\section{Discussion}

The failure is not an artefact of small samples. A crowd of experts estimates the population-average prediction, and a participant whose belief differs from that average is rewarded for moving the crowd's prediction toward its own belief. In case~A it does so by withholding its answer. More experts make that target more stable. Any repair has to bound the incentive by the heterogeneity of $\cD$, for example the divergence between individual priors and the crowd's average prediction, rather than by $m$ and $n$.

During the public review, one reviewer noted that the unweighted method seemed to contradict Theorem~2 in an experiment, and the authors attributed this to experts' and participants' priors not being drawn from the same distribution \citep{openreview2026forum}. Both counterexamples here satisfy that hypothesis.

\paragraph{What is not claimed.} Theorem~1 (common prior) is not affected. We do not evaluate the experiments of \citet{qiu2026truthfulness}, the practical usefulness of the method, or its claim about training dynamics. Our own verification is not an independent replication.

\paragraph{Acknowledgement.} Peer prediction for language models, its construction, Theorem~2 and the experiments are the work of Qiu, Carroll and Allen.

\begin{thebibliography}{10}
\bibitem[de~Moura and Ullrich(2021)]{demoura2021lean4} de~Moura, Leonardo, and Sebastian Ullrich. 2021. ``The Lean 4 Theorem Prover and Programming Language.'' In \emph{Automated Deduction -- CADE 28}, 625--635. Springer.
\bibitem[Goodhart(1975)]{goodhart1975problems} Goodhart, Charles A.~E. 1975. ``Problems of Monetary Management: The U.K. Experience.'' In \emph{Papers in Monetary Economics}, vol.~I. Reserve Bank of Australia.
\bibitem[Hoeffding(1963)]{hoeffding1963} Hoeffding, Wassily. 1963. ``Probability Inequalities for Sums of Bounded Random Variables.'' \emph{Journal of the American Statistical Association} 58(301):13--30.
\bibitem[mathlib Community(2020)]{mathlib2020} The mathlib Community. 2020. ``The Lean Mathematical Library.'' In \emph{Proceedings of the 9th ACM SIGPLAN International Conference on Certified Programs and Proofs (CPP)}, 367--381.
\bibitem[Miller et~al.(2005)]{miller2005peer} Miller, Nolan, Paul Resnick, and Richard Zeckhauser. 2005. ``Eliciting Informative Feedback: The Peer-Prediction Method.'' \emph{Management Science} 51(9):1359--1373.
\bibitem[OpenReview(2026)]{openreview2026forum} OpenReview. 2026. Discussion forum of ``Truthfulness Despite Weak Supervision'', ICLR 2026, submission 17499. \url{https://openreview.net/forum?id=mdw0vvRBEL}.
\bibitem[Qiu et~al.(2026)]{qiu2026truthfulness} Qiu, Tianyi Alex, Micah Carroll, and Cameron Allen. 2026. ``Truthfulness Despite Weak Supervision: Evaluating and Training LLMs Using Peer Prediction.'' In \emph{International Conference on Learning Representations (ICLR)}. arXiv:2601.20299.
\bibitem[Schoenebeck and Yu(2022)]{schoenebeck2022two} Schoenebeck, Grant, and Fang-Yi Yu. 2022. ``Two Strongly Truthful Mechanisms for Three Heterogeneous Agents Answering One Question.'' \emph{ACM Transactions on Economics and Computation} 10(4):1--26. doi:10.1145/3565560.
\bibitem[Strathern(1997)]{strathern1997improving} Strathern, Marilyn. 1997. ```Improving Ratings': Audit in the British University System.'' \emph{European Review} 5(3):305--321.
\end{thebibliography}

\end{document}
